% !PRINTED 417
\weberpage{417}
%<<CONT
eine Formel, die sich leicht durch vollständige Induction folgendermaassen
verallgemeinern lässt.

Bezeichnet man mit $\mu_1$ alle Zahlen, die aus $n$ entstehen, wenn man $n$
durch eine gerade Zahl verschiedener Primtheiler von $n$ dividirt ($n$ selbst
eingeschlossen), mit $\mu_2$ die Zahlen, die aus $n$ entstehen, wenn man $n$
durch eine ungerade Zahl solcher Primtheiler dividirt, so ist

\weq{10}{X_n = \dfrac{\prod\limits_{\mu_1}(x^{\mu_1} - 1)}{\prod\limits_{\mu_2}(x^{\mu_2} - 1)}.}

Der Beweis ergiebt sich aus (9) für $npp'$, wenn man annimmt, die Richtigkeit
sei für $n$ schon bewiesen.

Bedeutet $r$ jede der primitiven $n^{\mathrm{ten}}$ Einheitswurzeln, so sind
bei \gesperrt{ungeradem} $n$ die $-r$ die primitiven $2n^{\mathrm{ten}}$ Einheitswurzeln
(\S.\,132, I, II). Demnach ist bei ungeradem $n$

\weq{11}{X_{2n}(x) = X_n(-x).}

Ist $n$ eine Potenz von $2$, so ist nach (8):

\weq{12}{X_n = \dfrac{x^n - 1}{x^{\frac{n}{2}} - 1} = x^{\frac{n}{2}} + 1.}

Setzen wir in der Formel (8) $x = 1$, so erhält $X_n$ den Werth $p$, wenn $n$
eine Potenz von $p$ ist. Dagegen ergiebt die Formel (9), wenn $n > 1$ ist, für
$X_{npp'}$ den Werth $1$ (für $n = 1$ würde auf der rechten Seite Zähler und
Nenner $= 0$).

Wir erhalten also den Satz:

\gesperrt{V. Die Function $X_n$ erhält für $x = 1$ den Werth $p$, wenn $n$
eine Potenz der Primzahl $p$ ist, und den Werth $1$, wenn $n$ mehr als eine
Primzahl als Theiler hat.}

\begin{center}
\S.\,134.

\gesperrt{Irreducibilität.}
\end{center}

Die Functionen $X_n$, die im vorhergehenden Paragraphen definirt sind, haben
die für ihr tieferes Studium und alle Anwendungen sehr wichtige Eigenschaft
der \gesperrt{Irreducibilität}, d. h.

\gesperrt{Es ist nicht möglich, eine Function $X_n$ von $x$ in Factoren niedrigeren
Grades zu zerlegen, die selbst}
%CONT>>
